% Generated by scripts/assemble.py; corrections belong in source/corrections.json.
\examhead{181}{西北工业大学 2024 年数学分析试卷}{171}
% p172-b01; source PDF page 172
\par\Needspace{5\baselineskip}
\blockanchor{p172-b01}
\majorquestion{一}计算题.\par
% p172-b02; source PDF page 172
\blockanchor{p172-b02}
(1) 求极限 $\lim_{n\to \infty}\sum_{i = 1}^{n}\left(\frac{n + i}{n}\right)\sin \frac{i}{n^2}.$\par
% p172-b03; source PDF page 172
\blockanchor{p172-b03}
(2) 已知 $\lim_{x\to +\infty}\left(\frac{x - a}{x + a}\right)^x = \int_a^{+\infty}xe^{-2x}\mathrm{d}x,$ 求 $a$ 的值.\par
% p172-b04; source PDF page 172
\blockanchor{p172-b04}
(3) 计算积分 $\int_0^{\frac{\pi}{2}}\frac{1}{1 + \tan^{2024}x}\mathrm{d}x.$\par
% p172-b05; source PDF page 172
\par\Needspace{5\baselineskip}
\blockanchor{p172-b05}
\majorquestion{二}证明 $f(x) = \sin x^2$ 在 $(- \infty, +\infty)$ 上不一致连续.\par
% p172-b06; source PDF page 172
\par\Needspace{5\baselineskip}
\blockanchor{p172-b06}
\majorquestion{三}解答如下问题:\par
% p172-b07; source PDF page 172
\blockanchor{p172-b07}
(1) 设正项级数 $\sum_{n=1}^{\infty} a_n$ 收敛, $\{a_n\}$ 是单调数列, 证明: $\lim_{n \to \infty} na_n = 0$.\par
% p172-b08; source PDF page 172
\blockanchor{p172-b08}
(2) 已知 $A_{n} = \sum_{k=1}^{n} a_{k}$, 且 $\lim_{n \to \infty} A_{n}$ 存在, 证明: $\lim_{n \to \infty} \frac{a_{1} + 2a_{2} + \cdots + na_{n}}{n} = 0$.\par
% p172-b09; source PDF page 172
\par\Needspace{5\baselineskip}
\blockanchor{p172-b09}
\majorquestion{四}设 $f(x)$ 在 $[0,1]$ 上连续, 在 $(0,1)$ 上可导, 且 $f(0) = 0, f(1) = 1$.\par
% p172-b10; source PDF page 172
\blockanchor{p172-b10}
(1) 证明: 存在互异的 $\eta_1, \eta_2 \in (0,1)$, 满足 $\frac{1}{f'(\eta_1)} + \frac{1}{f'(\eta_2)} = 2$.\par
% p172-b11; source PDF page 172
\blockanchor{p172-b11}
(2) $a, b$ 为任意正实数, 证明: 存在互异的 $\xi_1, \xi_2 \in (0,1)$, 满足 $\frac{a}{f'(\xi_1)} + \frac{b}{f'(\xi_2)} = a + b$.\par
% p172-b12; source PDF page 172
\par\Needspace{5\baselineskip}
\blockanchor{p172-b12}
\majorquestion{五}解答如下问题.\par
% p172-b13; source PDF page 172
\blockanchor{p172-b13}
(1) 讨论级数 $\sum_{n=1}^{\infty} \frac{\sin n}{p^{\ln n}} (p > 0)$ 的敛散性.\par
% p172-b14; source PDF page 172
\blockanchor{p172-b14}
(2) 求幂级数  $\sum_{n=1}^{\infty}\frac{n^{2}}{(n+1)!}x^{n}$  的收敛域与和函数.\par
% p172-b15; source PDF page 172
\blockanchor{p172-b15}
(3) 设 $f$ 是以 $2\pi$ 为周期且在 $\mathbb{R}$ 二阶连续可导的函数, 证明: $f$ 的傅里叶级数在 $(- \infty, + \infty)$ 上一致收敛于 $f$.\par
% p172-b16; source PDF page 172
\par\Needspace{5\baselineskip}
\blockanchor{p172-b16}
\majorquestion{六}解答如下问题.\par
% p172-b17; source PDF page 172
\blockanchor{p172-b17}
(1) 讨论反常积分 $\int_0^{+\infty} \frac{\ln^m (1 + x)}{1 + x^n} \, \mathrm{d}x (n \geq 0)$ 的敛散性.\par
% p172-b18; source PDF page 172
\blockanchor{p172-b18}
(2) 求积分  $\int_{0}^{+\infty}\frac{\cos ax - \cos bx}{x}dx$ ，其中 0 < a < b.\par
% p172-b19; source PDF page 172
\blockanchor{p172-b19}
(3) 设 $F(z) = \iint_{x^2 + y^2 \leq z^2} f(x^2 + y^2) \, \mathrm{d}x \, \mathrm{d}y$, $f(0) = 1$, 且 $f$ 连续可导, 求 $F''(0)$.\par
% p172-b20; source PDF page 172
\par\Needspace{5\baselineskip}
\par\Needspace{6\baselineskip}
\blockanchor{p172-b20}
\majorquestion{七}设 $f(x) \in C(-\infty, +\infty)$, 定义函数列\par
\nopagebreak[4]
% p172-b21; source PDF page 172
\blockanchor{p172-b21}
\displayformula{f _ {n} (x) = \sum_ {k = 0} ^ {n - 1} \frac {1}{n} f \left(x + \frac {k}{n}\right), n = 1, 2, \dots}
% p172-b22; source PDF page 172
\blockanchor{p172-b22}
证明：函数列 $\{f_n(x)\}$ 在 $(- \infty, +\infty)$ 上内闭一致收敛于 $\int_0^1 f(x + t) \mathrm{d}t$.\par
